\section{Marco Teórico}

\subsection{Salario, reproducción social y condiciones materiales}

La investigación se inscribe en la economía política crítica y concibe el ingreso laboral como un recurso central para la reproducción material de las personas y los hogares. Desde la perspectiva marxista, el salario expresa una relación social de distribución y no únicamente un precio individual determinado por la productividad \parencite{MarxElCapital2008, DuranSilva2022}. Su suficiencia depende de las condiciones históricas de vida, de la organización del trabajo y del acceso a bienes y servicios necesarios.

Esta perspectiva permite situar la pregunta por el poder adquisitivo en la relación entre los ingresos monetarios y los recursos que hacen posible la vida cotidiana. La inflación, la segmentación del mercado laboral y la desigual distribución de los ingresos pueden modificar esa relación y trasladar a los hogares los costos de sostener su bienestar \parencite{DuranKremerman2024, cepalPanoramaSocialAmerica2025}.

\subsection{Canasta básica y capacidad de cobertura}

El costo de una canasta básica constituye una referencia institucional y estadística para estimar el valor monetario de necesidades definidas por cada país. No representa la totalidad del costo de vida ni una medida completa del bienestar: la composición, los requerimientos nutricionales, las escalas de equivalencia, la cobertura territorial y el procedimiento de actualización pueden variar entre fuentes nacionales.

En este estudio, la relación entre el ingreso monetario mensual per cápita y el costo mensual per cápita de la canasta local se utiliza como indicador descriptivo de capacidad de cobertura. La medida pregunta cuántas canastas definidas en cada contexto podrían cubrirse con el ingreso observado. No equipara las canastas entre países y no sustituye indicadores de pobreza multidimensional, acceso a servicios, riqueza o consumo efectivo.

\subsection{Comparación entre contextos nacionales}

Chile, Brasil y Argentina comparten desigualdades estructurales, pero presentan diferentes instituciones laborales, sistemas de protección social, trayectorias inflacionarias y metodologías estadísticas \parencite{cepalPanoramaSocialAmerica2025, CostaRibeiroCarcalhaes2024}. La comparación requiere conservar esas diferencias como parte del análisis, en vez de tratarlas como ruido que puede eliminarse mediante una conversión monetaria.

Por esta razón, los resultados se interpretan dentro del período, población y territorio cubiertos por cada encuesta y canasta. La evidencia disponible corresponde a cortes temporales no simultáneos; permite contrastar patrones descriptivos y discutir su contexto institucional, pero no identificar efectos causales ni una evolución común para los tres países.
